\documentclass[11pt]{article}

% ============================================================
% Response to Reviewer -- Round 18
% Companion to: revised_paper_v16.tex
% "Mortgage Lock-In and the Federal Reserve's Quantitative
%  Tightening Shortfall"
% Generated 2026-07-17 from the filled response letter in
% revision_roadmap_round18.md. Every number is drawn from a
% committed run artifact; section/line references are to
% revised_paper_v16.tex.
% ============================================================

\usepackage[margin=1in]{geometry}
\usepackage{amsmath}
\usepackage{amssymb}
\usepackage{enumitem}
\usepackage{parskip}
\usepackage[hidelinks]{hyperref}

% manuscript-label shorthand: renders tab:/fig:/eq: tokens in a
% lightweight typewriter form so they are clearly cross-references
% into the manuscript rather than this document.
\newcommand{\mlabel}[1]{\texttt{#1}}
\newcommand{\run}[1]{\texttt{#1}}

\title{\vspace{-2em}Response to Reviewer\\[0.3em]
\large Manuscript: \emph{Mortgage Lock-In and the Federal Reserve's
Quantitative Tightening Shortfall}\\[0.2em]
\normalsize Second-round referee report}
\author{}
\date{17 July 2026}

\begin{document}
\sloppy
\maketitle
\vspace{-2.5em}

Dear Editor and Reviewer,

Thank you for a careful and constructive report. Your summary of the paper's
central message --- that most of the QT MBS runoff shortfall was mechanical
and anticipated ex ante, with lock-in adding a robust but comparatively
modest $+9.2$ points --- is exactly the reading I intended, and several of
your requests sharpened exhibits that were already in the manuscript's
supporting layers but not yet surfaced as such. This revision adds eight new
computational exhibits and roughly fifteen manuscript edits; every new run
carries parity gates that reproduce the committed headline figures before any
parameter is varied. Below I respond point by point. Section and line
references (e.g.\ \S V.D, L522) are to the revised manuscript.

\section*{Technical weaknesses}

\subsection*{W1 --- $\beta_1=0$ switch-off and unobserved channels correlated with rate gaps}

The manuscript does not claim the $+9.2$pp marginal as causal identification,
and says so explicitly: Section~V.D's identification-status qualification
(L522) states the switch-off ``is not causal identification in the
econometric sense,'' that no quasi-experimental variation is exploited or
claimed, and that the marginal ``inherits any misspecification of the
held-fixed structure.'' Section~V.C (L448) goes further in exactly your
direction: because the design does not cancel a floor that co-varies with the
rate cycle, the identified marginal ``is best read as total rate-attributable
prepayment suppression, an upper bound on the strictly-voluntary elasticity's
contribution'' --- the channels you name load, if at all, \emph{into} that
upper bound rather than silently inflating a causal estimate. The Discussion
opener (L617) names servicer behavior and loan-level heterogeneity as
candidates, ``not demonstrated drivers of the level.'' Job-mobility shocks and
servicer-level strategies remain un-bounded candidates rather than controlled
channels; I state this plainly rather than claim otherwise. On the geographic
channel specifically, the marginal's regional structure is now measured
directly at the loan-sample level: the Census-region cells of run
\run{subgroup\_marginals} span per-balance intensities of only $0.93$--$1.06$,
so no region carries the marginal out of proportion to its balance.

\subsection*{W2 --- Ad-hoc freeze and DTI thresholds}

The ad-hoc status is conceded in the methodology section itself (\S III.C,
L144): the freeze's ``20\% share and 150bp trigger are author calibrations
with no external source,'' and the 43\% DTI is anchored to the pre-2021
General QM limit as ``a behavioral anchor rather than a live regulatory
constraint.'' Three quantitative layers now bound their influence.
(i)~Run \run{freeze\_sensitivity} (\S VII.C, L651): removing the freeze moves
ABM trapped liquidity \$84.5B $\to$ \$44.6B (isolated $+\$39.9$B/$+5.2$pp),
the 100/200bp trigger sweep spans \$105.9--\$70.5B, and re-anchoring the share
to the Liebersohn--Rothstein-implied 18.3\% external proxy moves recovery only
$-0.4$pp, with the involuntary floor bit-invariant throughout. (ii)~The DTI
wall itself is now swept 36/43/50\% on the same frozen surface
(run \run{dti\_threshold\_sweep}), moving ABM trapped liquidity to \$153.3B /
\$84.5B / \$119.7B (share $20.1/11.1/15.7\%$) with the involuntary floor
re-anchored into the empirical 4--5\% band at every threshold --- the gate is
material to the ABM's own estimate but never breaches the floor.
(iii)~Most importantly, the ABM is not the headline estimator (\S IV, L175):
the $+9.2$pp marginal is a hazard-framework object that contains no freeze and
no DTI gate.

\subsection*{W3 --- Freddie-only hazard data; overlay bounding ``not fully quantified''}

The bounding is quantified, in two signed layers printed in Section~V.C
(L450--452): a static composition bound of \$20--47B (2.6--6.2\% of benchmark,
signed toward overstating trapped liquidity), and a time-varying overlay that
scores the 20.4\% Ginnie face share by the observed Ginnie CPR series in both
Path~B legs --- the marginal is confined to the conventional share at $+7.3$
points, sign-positive under every variant tested (total-CPR, voluntary-only,
placebo), and the Ginnie-specific differential contribution is $+\$38.5$B,
5.0\% of benchmark, inside the static bound (see also \mlabel{tab:composition}'s
agency row and \S VII.F, L715). Cross-agency external replication on Fannie
Mae data reproduces the marginal at $+8.68$ points, inside the pre-committed
envelope (\mlabel{tab:specbox} Validation row, L338).

\subsection*{W4 --- Danish discounting: OAS sensitivity}

This sensitivity exists in the manuscript as Table~\mlabel{tab:discount}
(\S V.C, L456--470; run \run{danish\_discount\_bound}): discounting the Danish
buyback at primary minus 0/25/50/75/100bp flips $0/2{,}470/7{,}805/17{,}493/
32{,}488$ of $1{,}683{,}082$ loan-month PV classifications and moves trapped
liquidity by exactly \$0.00 billion on every accounting basis. The zero is
structural, not numerically small: the production Danish leg prepays at the
imported Berger-recalibrated flat elasticity, and no downstream quantity
consumes the PV classification, so ``option-adjusted marks $\ldots$ cannot
move the production $+\$61.2$ billion gap'' (printed at L456). At your
suggestion I extended the grid to the 150bp endpoint you name: primary minus
125bp flips $55{,}210$ and minus 150bp flips $92{,}495$ of the $1{,}683{,}082$
classifications, with trapped liquidity moved by exactly \$0.00B at both
points --- the structural zero holds across the full 0--150bp span, and
Table~\mlabel{tab:discount} now prints all seven rows. The OAS qualification
itself is stated where the proxy is introduced (L138), and full
option-adjusted duration/convexity is explicitly scoped out
(\mlabel{tab:wal} note, L742).

\section*{Experimental / methodological weaknesses}

\subsection*{M1 --- Subgroup reporting (coupon / vintage / LTV / FICO / geography)}

Coupon and vintage are decomposed at HEAD: the group-ablation grid (\S V.D
fourth qualification, L511; Figure~12) shows all fifteen
vintage~$\times$~coupon cells contribute positively, no cell exceeds 21\% of
the total, and the 2020--21 vintages carry 72\% of the marginal on 73\% of
balance (per-balance intensity $0.8$--$1.5\times$); the month~$\times$~coupon
split (Figure~11) adds terminal by-coupon margins of \$26.4B / \$30.4B /
\$12.7B ($<3.0\%$ / 3.0--4.0\% / $\geq 4.0\%$), positive in all 42 months. At
your request I extended the same engine to the remaining dimensions
(run \run{subgroup\_marginals}), and the outcome is the one you hoped would
bolster generality: every cell of the FICO ($<680$ / 680--740 / $740+$), LTV
($\leq 80$ / $>80$), and Census-region partitions contributes positively; the
per-grid cell sums reproduce the $+\$70.3$B total within interaction residuals
of 0.8, 0.4, and 0.5 percent (all below the vintage~$\times$~coupon grid's
1.0\%); and every cell's per-balance intensity lies in $0.88$--$1.26$, tighter
than the vintage~$\times$~coupon span. The largest single cell is the $740+$
FICO bucket at \$40.4B --- 57.9\% of the cell sum on 65.8\% of balance
(intensity $0.88$) --- and the mild tilts run toward below-740-FICO
($1.23$--$1.26$) and above-80-LTV ($1.23$) borrowers, with regional
intensities spanning only $0.93$--$1.06$. No segment carries the marginal out
of proportion to its balance; the manuscript now reports this alongside the
vintage~$\times$~coupon grid in Section~V.D.

\subsection*{M2 --- Single literature elasticity band}

Three clarifications and one addition. First, the paper does run a non-L\&R
elasticity through the production microsimulation: Berger (2026)'s
separately-estimated flat moving elasticity is the production Danish
counterfactual leg (L456; L791). Second, an in-sample elasticity estimated on
the same data already exists: Path~A's rate-gap coefficient, fit directly on
the Freddie panel, is $+0.65$ standardized ($+0.29$ per 100bp), positive in
197 of 198 bootstrap replications (\mlabel{tab:bootstrap}), with the Fannie
external refit reproducing the sign on a second agency at ${\sim}3.6\times$
the magnitude --- I now present this explicitly as sign-triangulation of the
L\&R band (\S V.B--V.C), with the caveats that a prepayment-hazard coefficient
is not unit-comparable to L\&R's quarterly-mobility decline and that Path~A
measures prepayment, not pure mobility. Third, Fonseca (2024)'s
${\sim}9\%$-per-100bp estimate is cited as independent corroboration of the
magnitude (footnote at \mlabel{eq:beta1}). Fourth --- executed at your
suggestion --- run \run{fonseca\_band\_anchor} propagates the Fonseca point
through the same transform and the same band machinery: it implies
$\beta_1 \approx 0.096$ and a marginal of $+\$87.9$B ($+11.5$ points), above
the L\&R high edge ($+\$79.5$B/$+10.4$ points) and inside the pre-committed
box, so the adopted central calibration is the conservative choice among the
available literature estimates; because that estimand is a moving-rate
reduction rather than a quarterly-mobility decline, I read it as directional
corroboration, not a competing point estimate. The FICO/LTV covariate priors
are also re-estimated in-sample ($\hat\beta_F = -0.39$, $\hat\beta_L = +0.21$),
with the simulation re-run under them (101.3\%) printed as the specification's
largest single sensitivity.

\subsection*{M3 --- Linear curtailment mapping}

The lock-in marginal is invariant to the mapping by construction, and this is
now demonstrated rather than asserted: run \run{curtailment\_profile\_demo}
(\S VII.F, L713) wraps the production curtailment series in a $\pm 25\%$
seasonal profile and a 2022--23 vs 2024--25 regime split --- the wedge moves
(9.16 and 8.90 vs 9.10 points; netted totals \$70.06B and \$68.09B) while the
marginal is unchanged to below $1.4\times10^{-14}$ points and the simulated
CPR paths are bit-identical. Any alternative functional form of the
income-to-rate map produces a monthly rate series, which is a special case of
the same identity: it can move only the common wedge, never the marginal or
any cross-estimator gap (now stated explicitly at \S III.D/VII.F). The level
itself is stressed $\pm 25\%$ by run \run{curtailment\_danish\_scaling}
(\S VII.F, L713), moving the netting \$52.2/\$69.6/\$87.0B with CPR paths
bit-identical.

\subsection*{M4 --- TBA/pool-homogeneity: quantitative deepening}

The liquidity-vs-mobility exchange is developed and cited (intro L45;
conclusion L793--795): Vickery (2013) documents that TBA-eligible pools trade
at significantly lower cost than comparable ineligible specified pools, and a
market-value repurchase rule would introduce path-dependent borrower-specific
payoff values that strain the fungibility this premium depends on. A formal
liquidation-option valuation under TBA deliverability rules would require the
term-structure and prepayment-option model this paper deliberately does not
build (\mlabel{tab:wal} note, L742); I cite the documented premium as the
relevant magnitude and leave the option-theoretic valuation as future work
rather than supply an unmodeled number.

\section*{Clarity}

\subsection*{C1 --- Hazard specifications and $\beta_1$ notation}

The full specifications are printed in the manuscript: Path~A's regression
with every term defined (\mlabel{eq:pathA}, L348--351, including the age-spline
basis and knots and the burnout orthogonalization); Path~B's cause-specific
hazards and the $\beta_1$ transform
(\mlabel{eq:pathB}/\mlabel{eq:default}/\mlabel{eq:beta1}, L435--445);
symbol-by-symbol notation with units (\mlabel{tab:hazard-notation}, L279;
$\beta_1$ ``log-hazard, per 100bp,'' tablenotes L313); the compact spec box
with the full FE set, dummies, and time-varying covariates
(\mlabel{tab:specbox}, L323); stratum construction at L347 (296 observed
vintage~$\times$~coupon~$\times$~FICO~$\times$~LTV cells). I have added one
sentence printing the numeric bin edges and stating explicitly that equity
enters only as static original LTV, with no local-HPA or dynamic-equity
covariate (\S V.B; see also Q1/Q9).

\subsection*{C2 --- Consolidated basis/inputs table}

Three exhibits carry this: \mlabel{tab:headline} (Section~I: every core result
with its uncertainty, cross-referenced to its home exhibit),
\mlabel{tab:bases} (\S V.D: the explicit four-basis crosswalk for the hazard
legs, with the \$69.56B netting stated), and the run ledger plus manifest
footnote (\mlabel{app:ledger}, \mlabel{fn:manifest}: inputs and provenance per
figure). I have closed the literal ask with a new crosswalk table
(\mlabel{tab:crosswalk}, in the run-ledger appendix) mapping each headline
quantity to its accounting basis and source run, with \mlabel{tab:headline}'s
caption pointing to it.

\subsection*{C3 --- Earlier empirical motivation for calibration choices}

The friction mapping and the freeze/DTI parameters are motivated --- and their
author-calibrated status disclosed --- in the methodology section (\S III.C,
L142--144), which precedes the ABM results (\S IV, L173). The
disclosure-of-limits placement is deliberate: I prefer a stated limitation
ahead of the results to a post-hoc defense behind them.

\section*{Related work}

\subsection*{L1 --- Modern ML survival baseline}

See R1 below and Section~V.B (L407): the flexible-learner ceiling test (run
\run{ml\_comparator\_holdout}) already quantifies the parametric headroom at
the paper's own unit of analysis --- a gradient-boosted Poisson learner
improves exposure-weighted RMSE to 0.93 against the frozen spec's 3.04 while
leaving unweighted RMSE (37.6 vs 37.9) and out-of-sample $R^2$ ($+0.003$)
unmoved; verdict: flexible fit helps, no headline change. A faithful
discrete-time-neural or MTLR baseline of the Sadhwani class trains on
loan-month histories (${\sim}826$M rows) that are external to this repository
under its data posture --- a decline of record I maintain, with the HGB result
as the tabular upper bound on what any flexible learner could change.

\subsection*{L2 --- MBS pricing/OAS and convexity hedging}

The connection is in the manuscript: the 2022--23 primary-to-10-year spread
widening (roughly 170 to over 300bp) with its decomposition attributed to
Boyarchenko (2019) (L47, with the printed caveat that the source predates the
2022--23 magnitudes), and the convexity-hedging feedback passage at L793
(Gabaix 2007; Malkhozov 2016; Perotti 2024). OAS pricing itself is scoped out
explicitly (L138, L742) and shown incapable of moving the Danish result
(\mlabel{tab:discount}). To situate the contribution within the broader
QT/balance-sheet-unwind literature, I have added two references and a situating
sentence to the transmission strand (L96): Krishnamurthy and Vissing-Jorgensen
(2011), whose channel decomposition of quantitative easing includes an
agency-MBS prepayment channel, and Vayanos and Vila (2021), whose
preferred-habitat model formalizes why the outstanding supply of held bonds
moves the term structure --- the lock-in friction I measure governs the pace at
which the Federal Reserve's MBS holdings run off, and so is an input to that
balance-sheet-supply channel rather than a phenomenon separate from it.

\subsection*{On the household lock-in / mobility literature}

Two references you suggest are already cited and integrated: Ferreira, Gyourko,
and Tracy (2010) at L80 (rising rates and negative equity each depress owner
mobility --- the equity-lock-in channel complementary to the rate channel I
model) and DeFusco and Mondragon (2020) at L86 (employment-documentation and
closing-cost frictions that inhibit refinancing pass-through, most severely in
recessions).

\subsection*{On future-dated sources}

The genuine preprints among the 2025--2026 entries now carry explicit status:
Peng and Lessmann (2026) and Lesniewski (2026) are marked ``Preprint,
arXiv:$\ldots$'' and Berger et al.\ (2026) ``Working paper, SSRN.'' The
remaining future-dated entries are published journal articles (Liebersohn 2024,
JFE 2025; Perotti 2024, JCAM 2026), left as such.

\subsection*{On reported precision and the magnitude of the marginal}

I agree that the magnitude is not tightly pinned, and the paper already takes
that position rather than the opposite: the narrow sampling interval
$[+9.17, +9.23]$ is explicitly labelled as sampling noise that is ``negligible
at reporting precision, so the operative uncertainty is the calibration bound''
(intro L43; Path~B L479), and the headline is stated as ``the sign and a
bounded range $\ldots$ not a pinned magnitude'' (abstract L30). The operative
uncertainty is the calibration box ($+2.1$ to $+13.2$ points), not the sampling
interval. This revision strengthens that framing with the round-18
regime-split result: because the baseline hazard multiplies the very channel
the null switches off, a $\pm 20\%$ structural break in the baseline moves the
marginal across $+5.0$ to $+12.1$ points, and I have added this
baseline-conditionality range to the consolidated uncertainty table
(\mlabel{tab:uncertainty}) alongside the calibration box, the floor-form value,
and the cyclical-floor range. The point value is conditional on the calibrated
baseline; the sign and the calibration-box membership are what the design
robustly supports.

\subsection*{On density and foregrounding (presentation)}

On front-loading the key definitions ($\beta_1$, the shared accounting basis,
the Path~A/B roles) and reducing parenthetical density: each is defined in the
manuscript (\S V.C \mlabel{eq:beta1} for $\beta_1$; \S III.D for the shared
accounting basis; \S IV--V for the two paths), and the non-binding-cap logic,
the Fed-projection justification, and the ABM-versus-hazard reconciliation are
all present and foregrounded in the abstract, introduction, and \S III.D
(L30, L43, L169) rather than buried. Where those points could be surfaced even
earlier, the trade-off is against the manuscript's deliberately compressed
style; I have flagged this as an editorial judgment rather than restructured the
exposition unilaterally, and welcome the editor's guidance on how much
front-loading the venue prefers.

\section*{Detailed comments}

\subsection*{T1 --- Confounders co-moving with rate levels; geographic/servicer controls}

See W1: the paper answers by withdrawing the causal claim rather than claiming
confounders are ruled out (L522, L805). Cross-agency composition stability on
geography and credit (State PSI 0.018, FICO 0.005, LTV 0.000;
\mlabel{tab:composition} panel~b) bounds selection drift; book-level
geographic/servicer reweighting is infeasible from CUSIP-level holdings (par
plus security description only) --- a decline of record. The geographic half
of the heterogeneity concern is now answered directly by the regional marginal
cells (see W1: intensities $0.93$--$1.06$), which measures the object the
requested Path~A regional fixed effects would only proxy.

\subsection*{T2 --- Stratum definitions and balance over time}

Definitions: L347 plus \mlabel{tab:specbox}; the numeric bucket edges are now
printed in Section~V.B. Balance: \mlabel{tab:panel} (train/holdout
stratum-month accounting, 10,176/6,077 cells with prepaid-UPB and exposure
totals) and \mlabel{tab:composition} (sample-vs-book and cross-agency PSI) are
the balance evidence; the grouped-calibration exhibit (Q10) adds the
per-half-year view of how the panel's fit behaves over time.

\subsection*{T3 --- Competing-risks equations and validation}

The explicit taxonomy is printed at \S V.C L446: Path~B is a discrete-time
multi-state model parameterized entirely through cause-specific hazards
(Putter 2007; implemented as in PyDTS, Meir 2025); no Fine--Gray
subdistribution hazard appears because every cumulative-incidence functional
is obtained by forward composition; the cause-specific reading is exact for
default and holds at the balance-aggregate level for prepayment (max hazard
sum 0.016 over 1.68M loan-months, so exact-exponential cause allocation
differs only at second order). On refi-vs-turnover: Path~B deliberately models
a single voluntary-prepayment hazard --- the rate-gap term carries the
refinance channel and the involuntary floor proxies turnover --- so a
refi/turnover cause-specific split is not identified in this design; the
external anchor I do have is the staff-measured scheduled-plus-turnover
decomposition (Na 2024, L92), which is the object the $\beta_1=0$ null
quantifies. The grouped-calibration exhibit now exists --- see Q10.

\subsection*{T4 --- Curtailment: seasonal/servicer sensitivity}

Executed: see M3. Seasonal and regime-split variants preserve the cancellation
to machine precision. Servicer heterogeneity is subsumed by the same identity:
curtailment enters as a single aggregate monthly rate on the shared holdings
path, netted identically from both legs, so any servicer-heterogeneous pattern
that aggregates to a monthly rate is a special case --- and a named
servicer-mix profile is now the demonstration's fourth profile: fixed
0.60/0.40 shares with front-loaded and flat group factors move the wedge to
9.30 points (\$71.12B netted) with the marginal again unchanged to
$1.4\times10^{-14}$ points.

\subsection*{E1 --- Out-of-time backtests}

What exists: the temporal month-holdout and isotonic recalibration on 2024+
months (\S V.B L407); the Fannie cross-agency replication ($+8.68$pp,
in-envelope; \mlabel{tab:specbox} Validation row L338); and --- closest to
your ask --- run \run{vintage\_residual\_bound}, which prices observed
QT-window speeds of cohorts outside the calibration universe: 2022 vintage
4.476\% vs sampled 4.303\% ($+0.173$pp) and pre-2017 5.264\% ($+0.962$pp), a
bound of \$11.75B (1.54\% of benchmark) signed toward overstating trapped
liquidity. I am candid that no untouched evaluation months exist anywhere in
the paper (\S VIII.A, L805). At your prompting I also isolated the 2015--16
cohorts specifically (run \run{vintage\_1516\_subleg}): they are 99.99\% of
the pre-2017 tail's exposure, prepaid at 5.26\% over the window ($+0.96$ points
against the sampled universe), and reproduce the tail's \$8.44B leg and the
\$11.75B bound essentially unchanged (the residual pre-2015 tail is 332
loan-months across 2010--2013) --- the cohorts you name sit inside the bound
the paper already prices.

\subsection*{E2 --- Subgroup decompositions}

See M1.

\subsection*{E3 --- ABM timing falsification}

The timing dimension carries four falsification exhibits, all printed: the
first-difference velocity check (simulated $\Delta$CPR correlates $-0.94$ with
contemporaneous rate changes; the empirical path's $+0.25$ is inside the
$\pm 0.31$ zero band --- the residual is not a matched velocity response,
L488); the phase-shift rerun (shifting the rate input $\pm 1$--$3$ months
never moves the peak cross-correlation to lag 0; L490 plus \mlabel{fig:ccf});
Theil $U_2$ on first differences ($1.226/1.010/1.003$ --- no estimator beats a
naive no-change forecast; \mlabel{tab:theil}); and the floor-anchor concession
you flag --- without it the ABM's CPR at 8\% market rate is 0.0\% and recovery
is 150.6\%, so the model ``earns no evidential credit'' on the floor dimension
and its exposure lies on the timing of the CPR path (L148). At your request I
extended this from the aggregate path to the cohort level (run
\run{cohort\_timing\_diagnostic}): across the six 30-year coupon cohorts
spanning the SOMA low-coupon mass (2.0--4.5\%), every simulated cohort CPR
tracks the contemporaneous rate change mechanically (first-difference
correlation $-0.87$ to $-0.94$; five of six cohorts phase-locked at lag 0,
the 4.5\% coupon peaking at lag $-6$ with a weaker $r=-0.29$, reproducing the
aggregate $-0.318$), while no realized cohort speed --- measured independently
in the Freddie 2017--2021 panel and the Fannie replication cohorts --- shows a
rate response distinguishable from zero at any lag within $\pm 6$ months (Freddie
lag-0 correlations $-0.05$ to $-0.16$, Fannie $-0.03$ to $-0.14$, all inside the
$\pm 2/\sqrt{n}$ band; the maximum magnitude over all cohorts and all lags is
0.34, roughly a third of the model's weakest cohort correlation, with no
coherent sign or lag). So the aggregate empirical
non-response is not a composition artifact and the simulated lead is not a
datable phase shift at the cohort level. One honest scope limit travels with
this: realized cohort-level monthly speeds exist only in the loan-level agency
panels that anchor Path~B, not in the SOMA holdings the aggregate empirical
CPR is back-out from, so the cohort split runs in the Path~B universe --- which
is itself why the estimator-level timing falsification stays at the SOMA
aggregate. All seven parity gates reproduce the committed
$-0.94/+0.25/-0.318$ anchors bit-exactly.

\subsection*{E4 --- Structural break in baseline hazards}

Executed, with an honest answer that corrects my own first instinct: run
\run{regime\_split\_marginal} re-runs the paired central and $\beta_1=0$ legs
under a $\pm 20\%$ multiplicative break in the baseline hazard from January
2023 onward (the post-spread-repricing boundary; the pandemic predates the
simulation window), imposed identically on both legs so nothing is fitted to
the benchmark. The marginal does not cancel: because the baseline multiplies
the rate-gap channel the null switches off, it scales with the baseline level
--- $+4.97$ points (\$37.97B) at $0.8\times$ to $+12.12$ points (\$92.72B) at
$1.2\times$, around the committed $+9.198$ points reproduced bit-exactly at
the unit factor --- while a $\pm 20\%$ break in the involuntary floor alone
moves it oppositely, to $+6.57$/$+10.41$ points. Every perturbation leaves the
marginal sign-positive and inside the pre-committed calibration box ($+2.1$ to
$+13.2$ points), so what survives a non-stationary baseline is the sign and
the box membership, consistent with the paper's standing posture that the
robust content is the marginal's sign, not its magnitude; the point value is
conditional on the calibrated baseline level, and Section~V.D now says so. The
Fannie replication ($+8.68$ points on an entirely different loan population and
baseline) and the landmark/isotonic holdout complement this as
different-baseline evidence.

\section*{Related-work comparisons}

\subsection*{R1 --- Discrete-time NN / MTLR baseline}

See L1.

\subsection*{R2 --- Lesniewski burnout-selection theory}

The tie is explicit at three points: the identity stated in Section~II (L88),
the design tie in Section~V.B (L375: the 295 strata cut cross-sectional
dispersion to within-cell variation; the demeaned surviving-stock regressor
proxies residual composition drift; the exposure-offset outcome is itself a
survival-weighted pool hazard), and the ABM burnout falsification read as the
identity's degenerate closed-population limit (L213). I deliberately do not
report an empirical within-stratum variance-evolution decomposition: the
estimated burnout coefficient is statistically zero under both bootstrap
schemes, and I hold it as a control, not a finding.

\subsection*{R3 --- Landmarking + isotonic recalibration variant}

Executed, at your earlier suggestion's next step: run
\run{landmark\_isotonic\_holdout} (\S V.B L407) fits an exposure-weighted
isotonic map on the 2023 inner-validation year only and applies it to the
frozen spec's 2024+ holdout predictions: RMSE improves to 2.53 against 3.04
--- short of the pre-committed $0.80\times$ materiality threshold; verdict: no
material change. Peng (2026) is cited in Section~II as demonstrated ``on the
same Freddie Mac performance data my hazard paths use.''

\section*{Broader impact (B1--B3)}

I thank the reviewer for the framing of the welfare balance and the
stress-test implication; the latter --- exposure to scheduled amortization and
baseline turnover mechanics as the primary runoff driver in high-rate regimes,
with lock-in elasticity an incremental amplifier --- is precisely the policy
content of the decomposition.

\section*{Questions}

\begin{description}[leftmargin=0pt, itemsep=0.6em]

\item[Q1 --- Exact hazard models, $\beta_1$ definition, covariates.]
Printed in full: see C1. Equity proxies and local HPA: equity enters only as
static original LTV ($\beta_L$ prior $+0.10$; panel estimate $+0.21$, CI
$[-0.59, +1.00]$); no local-HPA or dynamic-equity covariate exists, by design,
in this turnover-dominated deep-out-of-the-money window (now stated explicitly
in \S V.C).

\item[Q2 --- Strata definition; sensitivity to alternative stratifications.]
Definitions: L347, \mlabel{tab:specbox}, with numeric edges now printed. On
re-stratification sensitivity I decline the refit and give the reasons:
Path~A is pre-committed as aggregate corroboration excluded from headline
figures (\S V.B), so a finer-binned refit would harden a diagnostic exhibit no
headline consumes; the headline marginal's own subgroup structure is now
reported cell-by-cell over five dimensions --- vintage, coupon, FICO, LTV, and
region (M1/Q3) --- which is the object whose stability the stratification
question is after; and the coupon axis's signal content is tested directly by
the label-permutation ablation (real $z = 13.5$, exact one-sided
$p = 0.0099$).

\item[Q3 --- Subgroup decompositions of the $+9.2$pp.]
See M1: coupon/vintage at HEAD (all 15 cells positive, $\leq 21\%$ per cell,
72\%/73\% balance), and FICO/LTV/region now executed with all nine cells
positive at per-balance intensities $0.88$--$1.26$ --- stability across every
dimension asked, with the mild tilts (below-740 FICO, above-80 LTV) identified
and quantified.

\item[Q4 --- Ginnie handling; upper/lower bounds.]
Quantified: no reweighting toward Ginnie (stated, L450); static bound
\$20--47B signed toward overstating trapped liquidity; GMAR overlay confines
the marginal to the conventional share ($+7.3$pp, all variants sign-positive)
with the Ginnie-specific differential $+\$38.5$B inside the static bound. The
bound is one-signed because observed Ginnie speeds exceed conventional speeds
in every window month (window-mean differential 2.14pp vs snapshot 1.2--2.8pp)
--- an empirically one-sided ``upper/lower'' (now noted in \S V.C).

\item[Q5 --- In-sample elasticity to triangulate the band.]
It exists and is now framed as such: see M2 ($+0.65$ standardized / $+0.29$ per
100bp, 99.5\% bootstrap sign-stability, Fannie sign replication at
${\sim}3.6\times$; caveats on unit comparability and prepayment-vs-mobility
stated).

\item[Q6 --- Microdata for freeze/DTI; empirical-proxy sensitivity.]
No microdata or survey source supports the precise values --- the manuscript
says so verbatim (L144) rather than implying otherwise. Freeze half: the
18.3\% L\&R-implied external proxy and the trigger/share sweeps are printed
(\S VII.C). DTI half: now run --- 36/43/50\% moves ABM trapped to
\$153.3B/\$84.5B/\$119.7B (share $20.1/11.1/15.7\%$), with the involuntary
floor re-anchored into 4--5\% at every threshold (\S VII.C, run
\run{dti\_threshold\_sweep}). Realized-DTI-at-move distribution: no total-debt
microdata is available in this design; the front-end-only limitation is
disclosed (\S VIII.A).

\item[Q7 --- Curtailment variants preserve cancellation?]
Yes --- demonstrated, not asserted: see M3/T4. Seasonal and regime-split
variants move only the wedge; the marginal is invariant to
$1.4\times10^{-14}$ points with bit-identical CPR paths and bit-identical
central-vs-null netting.

\item[Q8 --- Danish OAS 50--150bp.]
0--100bp printed with structural \$0.00 deltas (\mlabel{tab:discount}); the
125/150bp rows are now run and printed --- flips $55{,}210$ and $92{,}495$ of
$1{,}683{,}082$, deltas exactly \$0.00 (see W4).

\item[Q9 --- CLTV dynamics and $\beta_1$ loading.]
Equity is static original LTV only; CLTV dynamics are unmodeled (no HPA feed
exists in the pipeline), and I acknowledge this as a limitation rather than
bound it indirectly: the static FICO/LTV ablation brackets the covariate
block's influence (108.3\% / 101.3\% --- the largest single sensitivity,
printed as such), and the identification-status qualification states the
marginal inherits any such misspecification. An HPA-updated dynamic-equity
hazard is future work requiring a new data ingest.

\item[Q10 --- Time-series calibration diagnostics.]
CPR-side diagnostics at HEAD: \mlabel{tab:theil} ($U_1$ levels
$0.412/0.431/0.265$; $U_2$ first-differences $1.226/1.010/1.003$; MSE
decomposition; cumulative dollar-error profiles) and the grouped
train/post-training split (98\% vs 137\%, \S V.B L407). The requested grouped
exhibit is now run \run{grouped\_calibration} (reported in the dynamic-fit
appendix): frozen Path~A predictions binned into exposure-weighted prediction
quintiles (edges fit on the training window only) crossed with calendar
half-years --- the training halves calibrate at a pooled predicted-to-realized
ratio of 1.000 (per-half 0.80--1.17), the post-2024 holdout drifts
monotonically (0.64, 0.57, 0.44, 0.43; pooled predicted 2.24 vs realized 4.21
points), and the reliability slope is intact away from the drift (upper
prediction bins 0.94--1.11 pooled), localizing the $U_2 > 1$ aggregate
diagnosis to prediction level $\times$ half-year. As a Path~A diagnostic it
moves no headline. Default-side cumulative-incidence calibration is declined
with reasons: default enters only through the delinquency pipeline and the
benchmark object is prepay-driven runoff, so realized default paths are not a
headline quantity in this design.

\end{description}

\vspace{1em}

Every new exhibit above passed its parity gates against the committed headline
figures before any parameter was varied, and the headline decomposition ---
the 88.7\% mechanical null and the $+9.2$-point lock-in marginal --- is
unchanged by the additions. Where a new run sharpened the paper's honesty
rather than its headline (the regime-split result of E4 in particular, which
shows the marginal's point value to be conditional on the baseline level while
its sign and calibration-box membership are robust), I have said so in both
the manuscript and this letter. I am grateful for a report that improved the
paper on exactly the dimensions it asked about.

\vspace{1.5em}
\noindent Sincerely,\\[0.5em]
The Author

\end{document}
